% !TeX root = main.tex
\appendix
\renewcommand{\chaptername}{Додаток}
\chapter{Список публікацій здобувача за темою дисертації та відомості про апробацію результатів}
\label{app:publications_and_approbation}

\textit{Наукові праці, в яких опубліковано основні наукові результати дисертації}

\begin{enumerate}[label=\arabic*.]
\item Денисюк Д., Савенко Б., Каштальян А., Іванченко О. Метод виявлення зловмисного програмного забезпечення на основі аналізу поведінкових патернів системи. \textit{Вчені записки ТНУ імені В.~І.~Вернадського. Серія: Технічні науки}. 2024. Т.~35(74), вип.~5. С.~124--129. DOI: \url{https://doi.org/10.32782/2663-5941/2024.5.1/19}.
\item Денисюк Д., Сорочинський О., Гнатчук Є., Дрозд А. Інформаційна технологія виявлення зловмисних кодів в інформаційних системах на основі аналізу паралельних процесів. \textit{Вісник Хмельницького національного університету. Серія: Технічні науки}. 2025. Т.~347, №~1. С.~576--583. DOI: \url{https://doi.org/10.31891/2307-5732-2025-347-80}.
\end{enumerate}

\textit{Праці, які засвідчують апробацію матеріалів дисертації}

\begin{enumerate}[label=\arabic*.,start=3]
\item Denysiuk D., Savenko O., Lysenko S., Savenko B., Kashtalian A. Method for detecting steganographic changes in images using machine learning. \textit{Proceedings of the 2023 IEEE 13th International Conference on Dependable Systems, Services and Technologies (DESSERT-2023)}. 2023. P.~1--6. DOI: \url{https://doi.org/10.1109/DESSERT61349.2023.10416453}.
\item Denysiuk D., Savenko B., Kashtalian A., Savenko O., Lysenko S. Detection system of software implants using decoys. \textit{Proceedings of the 2024 IEEE 14th International Conference on Dependable Systems, Services and Technologies (DESSERT)}. Athens, Greece, 2024. P.~1--6. DOI: \url{https://doi.org/10.1109/DESSERT65323.2024.11122205}.
\item Denysiuk D., Sochor T., Kapustian M., Kashtalian A., Savenko O. Methods for detecting software implants in corporate networks. \textit{CEUR Workshop Proceedings}. 2024. Vol.~3675. P.~270--284. URL: \url{https://ceur-ws.org/Vol-3675/paper20.pdf}.
\item Denysiuk D., Sochor T., Kapustian M., Kashtalian A., Drozd A. A method for detecting botnets in IT infrastructure using a neural network. \textit{CEUR Workshop Proceedings}. 2024. Vol.~3736. P.~282--292. URL: \url{https://ceur-ws.org/Vol-3736/paper21.pdf}.
\item Denysiuk D., Savenko O., Lysenko S., Savenko B., Nicheporuk A. Detecting software implants using system decoys. \textit{CEUR Workshop Proceedings}. 2025. Vol.~3899. P.~277--287. URL: \url{https://ceur-ws.org/Vol-3899/paper25.pdf}.
\item Denysiuk D., Savenko O., Kashtalian A. Методи виявлення вторгнення в IT інфраструктуру. \textit{ITSec-2024}. 2024. С.~86--87. URL: \url{https://itsec-conf.org/ua/handle/17.itsec/article.2024a6e}.
\item Denysiuk D., Savenko O., Kvassay M. Method for detecting malicious commands transmitted via images using steganography. \textit{CEUR Workshop Proceedings}. 2025. Vol.~3963. P.~340--350. URL: \url{https://ceur-ws.org/Vol-3963/paper27.pdf}.
\item Denysiuk D., Savenko B. A study of methods for detecting hidden threats in multimedia objects on web resources. \textit{Proceedings of the 15th International Scientific Conference ``Information Technology Security'' (ITSec-2026)}. 2026. P.~175--177. Handle: \url{https://itsec-conf.org/ua/handle/17.itsec/article.20260ea}.
\end{enumerate}


\section*{Відповідність положень новизни публікаціям}
\phantomsection
\addcontentsline{toc}{section}{Відповідність положень новизни публікаціям}

Матрицю складено за повними текстами праць №~1, 2, 5--7, 9 і 10 та за офіційними бібліографічними записами праць №~3, 4 і 8. У ній розмежовано пряме змістове покриття, тематичну передумову та відсутність тотожної публікаційної опори. Авторство не використано як заміну експериментальному доказу, а тематичну близькість --- як підставу для приписування непідтверджених результатів.

\begin{landscape}
\footnotesize
\begin{longtable}{@{}p{0.16\linewidth}p{0.27\linewidth}p{0.11\linewidth}p{0.25\linewidth}p{0.14\linewidth}@{}}
\caption{Документальна відповідність положень наукової новизни публікаціям і доказовим матеріалам}\label{tab:app_novelty_publications}\\
\toprule
Положення новизни & Публікаційна опора і сторінки & Місце в роботі & Експериментальна опора & Доказовий висновок \\
\midrule
\endfirsthead
\multicolumn{5}{l}{Продовження таблиці \thetable}\\
\toprule
Положення новизни & Публікаційна опора і сторінки & Місце в роботі & Експериментальна опора & Доказовий висновок \\
\midrule
\endhead
\bottomrule
\endfoot
Модель мультимодального представлення прихованої загрози & №~10, с.~175--177: об’єкт $M=\{P,F,S,Meta,B\}$ і зважена оцінка ризику; №~3 і 9 --- статичні та поведінкові передумови & Розділ~\ref{chap:models_formalization} & Формальну узгодженість часткового представлення перевірено для статичних складників; повне $Z$ із причинно пов’язаними подіями не формувалося & Прямо опубліковано концептуальну основу; типізоване $Z$, походження, маска доступності й граничні стани розвинено в дисертації та не подано як емпірично завершену модель \\
Метод аналізу спектральних, зорових і структурних ознак & №~3, с.~1--6: виявлення стеганографічних змін; №~9, с.~340--350: частотний і поведінковий аналіз; №~10, с.~175--177: групи статистичних, структурних і частотних ознак & Підрозділ~\ref{sec:svs_method}, розділ~\ref{chap:experiments} & Скорочені числові проєкції перевірено внутрішньо й у чотирьох зовнішніх серіях; локалізації та повний $H$ не оцінено & Публікації підтверджують предметні передумови; дисертаційна двогілкова процедура має часткову емпіричну опору в установлених форматних межах \\
Метод поведінкового узгодження ознак різного походження & №~1, с.~124--129; №~2, с.~576--583; №~9, с.~340--350: поєднання статичних і поведінкових спостережень; №~10, с.~175--177: поведінкова складова мультимодального ризику & Підрозділ~\ref{sec:behavior_alignment}, розділ~\ref{chap:experiments} & Перевірено лише скорочену модель узгодження статичних оцінок; причинно пов’язаних журналів, повного $X_b$ і $Z_b$ немає & Опубліковано предметну передумову, але не тотожний причинно-фазовий метод; тому новизну заявлено як формалізацію, а кількісну перевагу повної конфігурації не заявлено \\
П’ятиетапна інформаційна технологія & №~2, с.~576--583: статичний і динамічний етапи технології; №~10, с.~175--177: комплексне оцінювання мультимедійного об’єкта & Підрозділ~\ref{sec:information_technology}, розділ~\ref{chap:experiments} & Локально перевірено переходи $S_1$--$S_5$, атомарну реєстрацію, повтор і відновлення; промисловий контур не перевірено & Публікації підтверджують технологічну передумову; розвиток у дисертації стосується типізованих станів, реєстрації до зовнішньої дії та нового циклу реконструкції \\
\end{longtable}

\begin{longtable}{@{}p{0.06\linewidth}p{0.25\linewidth}p{0.24\linewidth}p{0.36\linewidth}@{}}
\caption{Документально підтверджені відомості про публікації та особистий внесок}\label{tab:app_publication_contribution}\\
\toprule
№ & Змістовий зв’язок із дисертацією & Публічно підтверджені відомості про внесок & Межа використання у дисертації \\
\midrule
\endfirsthead
\multicolumn{4}{l}{Продовження таблиці \thetable}\\
\toprule
№ & Змістовий зв’язок із дисертацією & Публічно підтверджені відомості про внесок & Межа використання у дисертації \\
\midrule
\endhead
\bottomrule
\endfoot
1 & Поведінкові шаблони системи та два окремі механізми виявлення & Д.~О.~Денисюк --- перший автор; розподіл функціональних ролей у відкритому тексті не оприлюднено & Предметна передумова поведінкового аналізу; не доказ причинно-фазового узгодження \\
2 & Інформаційна технологія зі статичним і динамічним аналізом & Д.~О.~Денисюк --- перший автор; розподіл функціональних ролей у відкритому тексті не оприлюднено & Пряма технологічна передумова; не тотожний п’ятиетапний контракт станів $S_1$--$S_5$ \\
3 & Виявлення стеганографічних змін у зображеннях засобами машинного навчання & Д.~О.~Денисюк --- перший автор; офіційний бібліографічний запис не містить розподілу ролей & Пряма передумова статичного аналізу; не підтвердження повного $H$ або відеоформатів \\
4 & Виявлення програмних імплантів із використанням системних приманок & Д.~О.~Денисюк --- перший автор; офіційний бібліографічний запис не містить розподілу ролей & Тематичний контекст реакції середовища; не прямий мультимедійний результат \\
5 & Поведінкове виявлення програмних імплантів у корпоративних мережах & Д.~О.~Денисюк --- перший автор; у повному тексті зазначено рівний внесок усіх авторів & Тематична передумова поведінкового каналу; не перевірка другого методу дисертації \\
6 & Нейромережеве виявлення ботмережевої активності & Д.~О.~Денисюк --- перший автор; у повному тексті зазначено рівний внесок усіх авторів & Ширший контекст аналізу поведінкових послідовностей; не прямий мультимедійний результат \\
7 & Виявлення програмних імплантів за реакцією системних приманок & Д.~О.~Денисюк --- перший автор; у повному тексті зазначено рівний внесок усіх авторів & Тематична передумова контрольованого середовища; не доказ причинної прив’язки до мультимедійного об’єкта \\
8 & Методи виявлення вторгнень в інформаційній інфраструктурі & Д.~О.~Денисюк --- перший автор; у матеріалах тез розподіл функціональних ролей не оприлюднено & Оглядовий контекст; не пряме підтвердження жодного з чотирьох положень новизни \\
9 & Частотне виявлення прихованих команд і спостереження за поведінкою системи & Д.~О.~Денисюк --- перший автор; у повному тексті зазначено рівний внесок усіх авторів & Пряма передумова першого методу та часткова передумова другого; не тотожний причинно-фазовий контракт \\
10 & Піксельне, частотне, структурне, службове й поведінкове подання мультимедійного об’єкта & Д.~О.~Денисюк --- перший автор; у матеріалах конференції розподіл функціональних ролей не оприлюднено & Пряма опублікована концептуальна передумова моделі та комплексного ризику; не емпіричне підтвердження повного $Z$ \\
\end{longtable}
\end{landscape}

Таблиця~\ref{tab:app_publication_contribution} завершує документальне зіставлення на рівні відомостей, доступних у публічних першоджерелах. Вона не перетворює порядок авторів на доказ виконання конкретної операції. Підписана довідка про розподіл постановки задачі, формалізації, програмної реалізації, експерименту, статистичного опрацювання й підготовки тексту потрібна лише тоді, коли її окремо вимагає процедура закладу; без погоджених відомостей такий розподіл у дисертації не вигадується.

\section*{Відомості про апробацію результатів дисертації}
\phantomsection
\addcontentsline{toc}{section}{Відомості про апробацію результатів дисертації}

У наведеній нижче таблиці подано відомості, які підтверджуються опублікованими матеріалами відповідних наукових заходів. Позначення форми апробації не засвідчує усної або стендової доповіді, якщо це окремо не підтверджено програмою заходу чи документами здобувача.

\small
\begin{longtable}{@{}>{\centering\arraybackslash}p{0.08\textwidth}p{0.35\textwidth}p{0.24\textwidth}p{0.20\textwidth}@{}}
\caption{Відомості про апробацію опублікованих матеріалів}\label{tab:app_approbation}\\
\toprule
№ публ. & Науковий захід & Місце і дата проведення & Документально підтверджена форма оприлюднення \\
\midrule
\endfirsthead
\multicolumn{4}{l}{Продовження таблиці \thetable}\\
\toprule
№ публ. & Науковий захід & Місце і дата проведення & Документально підтверджена форма оприлюднення \\
\midrule
\endhead
\bottomrule
\endfoot
3 & 13-та Міжнародна конференція IEEE з надійних систем, сервісів і технологій DESSERT~2023 & Афіни, Греція; 13--15 жовтня 2023~р. & Опубліковано повний текст доповіді \\
4 & 14-та Міжнародна конференція IEEE з надійних систем, сервісів і технологій DESSERT~2024 & Афіни, Греція; 11--13 жовтня 2024~р. & Опубліковано повний текст доповіді \\
5 & 5-й Міжнародний семінар Intelligent Information Technologies \& Systems of Information Security (IntelITSIS~2024) & Хмельницький, Україна; 28 березня 2024~р. & Опубліковано повний текст доповіді \\
6 & 1-й Міжнародний семінар Intelligent \& CyberPhysical Systems (ICyberPhyS~2024) & Хмельницький, Україна; 28 червня 2024~р. & Опубліковано повний текст доповіді \\
7 & 1-й Міжнародний семінар Advanced Applied Information Technologies (AdvAIT~2024) & Хмельницький, Україна; Жиліна, Словаччина; 5 грудня 2024~р. & Опубліковано повний текст доповіді \\
8 & 13-та Міжнародна наукова конференція «Безпека інформаційних технологій» (ITSec~2024) & Львів, Україна; 9--11 травня 2024~р. & Опубліковано тези доповіді \\
9 & 6-й Міжнародний семінар Intelligent Information Technologies \& Systems of Information Security (IntelITSIS~2025) & Хмельницький, Україна; 4 квітня 2025~р. & Опубліковано повний текст доповіді \\
10 & 15-та Міжнародна наукова конференція «Безпека інформаційних технологій» (ITSec~2026) & 27--29 травня 2026~р. & Опубліковано тези доповіді, с.~175--177 \\
\end{longtable}

\clearpage
